\documentclass[12pt]{article}

% Page setup
\usepackage[margin=1in]{geometry}
\usepackage{setspace}
\usepackage{enumitem}
\usepackage{amsmath, amssymb}

% No paragraph indent
\setlength{\parindent}{0pt}

% Line spacing
\renewcommand{\baselinestretch}{1.2}

\begin{document}

\begin{center}
\textbf{\Large MATH 2790 Fall 2025}\\[1em]
\textbf{\Large Homework 13}\\[0.5em]
Due date: Wednesday, Dec 10
\end{center}

\bigskip

\textbf{Instructions:}

\begin{enumerate}[label=(\alph*)]
\item It is expected to show detailed work for every problem. You are allowed to have a
discussion with peers on some challenging questions. But then you should work on
your own to finish the assignment.

\item All questions in this assignment are from sections 8.3 and 9.2.
\end{enumerate}

\bigskip

\textbf{Problems:}

\begin{enumerate}
\item Question (2) on page 435 of textbook under section 8.3.

\item Question (3) on page 435 of textbook under section 8.3.

\item Question (5) on page 435 of textbook under section 8.3.

\item Question (2) on page 496 of textbook under section 9.2.

\item Question (3) on page 496 of textbook under section 9.2.

\item Question (4) on page 496 of textbook under section 9.2. It should be a par-value bond.

\item On December 31, year $X$, your company estimates that it will pay a total of \$10 million in year $X+1$ and subsequent years (from year $X+2$ to year $X+4$) on health insurance claims that were incurred in year $X$. The payout pattern is as follows:

\[
\begin{array}{c|cccc}
\text{Year} & X+1 & X+2 & X+3 & X+4 \\ \hline
\text{Payout percentage} & 60\% & 20\% & 10\% & 10\%
\end{array}
\]

Determine the Macaulay duration of the claim payments assuming that they are paid
at the end of each calendar year and the effective rate of interest is $10\%$.
\end{enumerate}

\end{document}
